\documentclass[UTF8,a4paper,12pt]{ctexart}

\usepackage{geometry}
\usepackage{graphicx}
\usepackage{booktabs}
\usepackage{longtable}
\usepackage{multirow}
\usepackage{array}
\usepackage{tabularx}
\usepackage{hyperref}
\usepackage{enumitem}
\usepackage{amsmath}
\usepackage{titlesec}
\usepackage{caption}
\usepackage[numbers,sort&compress]{natbib}
\usepackage{fancyhdr}
\usepackage{setspace}
\usepackage{tikz}
\usepackage{algpseudocode}

\usetikzlibrary{positioning, arrows.meta, fit, backgrounds}
\geometry{margin=2.4cm}
\hypersetup{colorlinks=true, linkcolor=black, citecolor=black, urlcolor=blue}
\setlength{\parindent}{2em}
\setlength{\parskip}{0pt}
\setlength{\emergencystretch}{3em}
\setlist{nosep}
\captionsetup{font=small, labelfont=bf}
\renewcommand{\bibname}{参考文献}
\renewcommand{\refname}{参考文献}
\renewcommand{\arraystretch}{1.15}

% -- Float placement: keep figures/tables with surrounding text --
\renewcommand{\topfraction}{0.85}
\renewcommand{\bottomfraction}{0.75}
\renewcommand{\textfraction}{0.10}
\renewcommand{\floatpagefraction}{0.80}
\setcounter{topnumber}{3}
\setcounter{bottomnumber}{2}
\setcounter{totalnumber}{5}
\setlength{\intextsep}{8pt plus 2pt minus 2pt}
\setlength{\floatsep}{10pt plus 2pt minus 2pt}
\setlength{\textfloatsep}{12pt plus 2pt minus 4pt}

\newcolumntype{L}[1]{>{\raggedright\arraybackslash}p{#1}}
\newcolumntype{Y}{>{\raggedright\arraybackslash}X}
\usetikzlibrary{arrows.meta,calc,fit,positioning,shapes.geometric}
\newcounter{algorithm}
\renewcommand{\thealgorithm}{\arabic{algorithm}}

% -- CCCF style: Chinese numeral section headings (一、二、三...) --
\renewcommand{\thesection}{\chinese{section}}
\titleformat{\section}{\centering\heiti\zihao{-3}}{\thesection、}{0em}{}
\titleformat{\subsection}{\heiti\zihao{4}}{\arabic{section}.\arabic{subsection}}{0.6em}{}
\titleformat{\subsubsection}{\heiti\zihao{-4}}{\arabic{section}.\arabic{subsection}.\arabic{subsubsection}}{0.5em}{}

\titlespacing*{\section}{0pt}{1.2em}{0.8em}
\titlespacing*{\subsection}{0pt}{0.8em}{0.4em}
\titlespacing*{\subsubsection}{0pt}{0.6em}{0.3em}

\pagestyle{fancy}
\fancyhf{}
\fancyfoot[C]{\thepage}
\renewcommand{\headrulewidth}{0pt}

\begin{document}

% ===========================================
% Title block (following CCCF 2601 reference style)
% ===========================================
\begin{center}
{\bfseries\zihao{2} 国产算力平台大模型推理引擎综述\par}
\vspace{0.95em}
{\zihao{4} 张书豪\textsuperscript{1}，廖小飞\textsuperscript{1}\par}
\vspace{0.45em}
{\zihao{5} \textsuperscript{1}大数据技术与系统国家地方联合工程研究中心\par}
{\zihao{5} 服务计算技术与系统教育部重点实验室\par}
{\zihao{5} 集群与网格计算湖北省重点实验室\par}
{\zihao{5} 华中科技大学计算机科学与技术学院\par}
\end{center}

\renewcommand{\thefootnote}{\fnsymbol{footnote}}
\footnotetext[1]{受新一代人工智能国家科技重大专项资助（课题编号：2025ZD0123804）\\ 通信作者：廖小飞，E-mail：xfliao@hust.edu.cn}
\renewcommand{\thefootnote}{\arabic{footnote}}

\vspace{0.6em}

\begin{spacing}{1.2}

% ===========================================
% Body (no abstract at top — CCCF style puts abstract at end)
% ===========================================
\section{引言}

大模型应用的快速扩张，正把推理系统从单机实验环境推向面向生产的高并发服务平台。随着国产 AI 加速芯片、服务器整机与集群软件生态逐步成熟，围绕国产算力构建高性能、可维护、可扩展的推理引擎，已成为学界和产业界共同关注的方向。但“国产算力推理引擎”并不是把既有 serving 框架迁移到另一类设备上的简单重复。相较通用硬件环境，国产平台往往同时叠加设备能力离散、编译链约束差异、运行时治理复杂和交付路径多样等因素，使许多原本可以局部吸收的问题迅速上升为系统级问题。平台边界、运行时状态与工程闭环一旦同时变化，推理系统的设计重心也会随之迁移。与训练系统相比，推理引擎更关心动态批处理、显存复用、流式时延、长上下文，以及工具调用、多模态和结构化输出下的稳定性；对国产平台而言，还要额外处理硬件后端差异、算子完备性、编译运行时兼容性，以及上游框架快速演进带来的维护压力。

近年来，以 vLLM 为代表的推理框架通过 PagedAttention、连续批处理与高效内存管理显著改写了大模型服务系统的设计范式，SGLang 等系统则进一步强调请求编排、结构化控制流和复杂工作负载下的执行效率\citep{kwon2023efficient,vllm,sglang}。通用 serving 文献也表明，推理系统的研究重心正在持续外移：Orca 与 Sarathi-Serve 说明，迭代级调度、连续批处理和 chunked prefill 已成为高吞吐服务的基础组织方式\citep{orca2023,sarathi2023}；DistServe 与 Splitwise 表明，prefill 和 decode 的阶段解耦正在从局部优化上升为整体 goodput 设计问题\citep{zhong2024distserve,splitwise2024}；XGrammar 与 ElasticMM 则进一步把演进推向结构化执行与多模态并行，说明推理引擎正在从 token 生成器演变为状态密集型运行时\citep{dong2024xgrammar,liu2025elasticmm}。这意味着，当推理系统进入国产芯片、国产编译运行时和本土部署环境后，关键问题会迅速从单一硬件后端优化转向统一框架抽象与本土设备能力的持续对齐。此时，调度、缓存、多模态和结构化执行的可迁移性、可维护性与稳定性，往往比单点优化更重要。从产业侧看，国产算力推理引擎面向的也早已不只是学术基准测试，而是在线问答、代码助手、政企知识服务、具身智能中间层和 Agent 工作流等多种 AGI4S 场景；这些场景既要求较高的吞吐密度和资源利用率，也要求流式首 token 时延、上下文切换开销、错误恢复能力和服务可观测性达到生产标准。因此，国产推理引擎的竞争焦点正从“能否完成基础部署”转向“能否在真实业务下稳定、高效且可持续演进地运行”。

在这一背景下，开源主干复用增强路线、平台一体化方案和国产原生/独立路线虽然都在解决国产部署问题，但三者分化的根源并不只是优化手段不同，更在于对系统边界和演进节奏的假设不同：前者强调上游复用，并允许在共享主干外侧持续注入平台与场景增强，其中既包括直接围绕主线后端边界展开的适配，也包括以 vLLM-HUST 为代表的本土增强实践；中者强调平台统一性交付；后者则更侧重围绕自研状态机、缓存治理和部署闭环组织系统能力。相较训练侧更容易用 FLOPS、集群规模等指标描述阶段进展，推理侧的核心矛盾更多分布在请求形态、缓存复用、阶段干扰、故障恢复与升级节奏等难以被单一数字概括的维度上。尤其当模型族从纯文本扩展到多模态、语音和 Agent 工作流后，系统能否在较长时间尺度内维持接口兼容、性能稳定和运维可控，往往比一次局部 benchmark 的领先更能决定其工程价值。基于此，全文采用“一个底座、三条主路径”的分析骨架：底座是统一抽象、统一指标和可观测证据链，用来支撑多平台、多版本和多阶段优化条件下的可组合、可比较与可复现；三条主路径分别是执行与通信、状态治理、压缩与成本协同，对应执行骨架收敛、长上下文与高并发下的状态组织，以及单位 token 成本进入系统闭环。在此基础上，后文依次讨论生态格局与路线分化、分析骨架与核心架构、基于骨架的路线比较、由路线比较收束出的关键挑战，以及未来可能的收敛方向。全文重点不在于为某一路线背书，而在于识别哪些能力应沉淀为能力契约，哪些状态必须纳入运行时统一管理，哪些优化适合以插件或外围工具链独立演进，以及哪些 benchmark 足以支撑路线选择。

为使讨论建立在可交叉核验的材料之上，本文还对讨论范围和证据边界作如下限定。首先，本文的“国产算力推理引擎”既包括直接面向国产芯片的原生推理系统，也包括基于主流 serving 主干的国产后端适配、面向多类国产硬件的部署工具链，以及围绕服务状态机、缓存治理和交付闭环组织的独立运行时；但不把单纯的硬件整机、一体机交付页面或应用层封装直接等同于推理引擎本体。其次，本文优先使用四类证据：系统论文与技术报告用于界定机制与设计目标，官方文档与开源仓库用于确认能力边界和工程入口，标准或白皮书用于辅助描述生态收敛趋势，社区部署记录仅作为观察真实摩擦点的补充材料，而不单独支撑性能结论。再次，文中对路线比较主要围绕可演进性、状态治理能力、复杂负载承载能力和交付确定性展开，而不试图在缺少统一复现实验条件的前提下给出跨平台绝对性能排序。本文更关注复杂度被压在哪里、证据是否足以支撑相关判断，以及系统是否具备可持续演进边界，而不是把不同公开口径简单折算为单一优劣结论。

\subsection{研究对象与证据方法}

本文的综述对象，限定为近年在公开论文、技术报告、官方文档或开源仓库中能够形成交叉验证的国产算力推理相关系统与工程实践。就时间范围而言，文章重点覆盖大模型 serving 主干与国产后端适配快速演进的近年公开材料，并在必要处回溯少量更早的基础工作，用于交代连续批处理、PagedAttention、阶段解耦和结构化执行等机制来源。就对象纳入而言，只有当某一对象至少能够在“机制说明”“工程入口”“公开实现”三者中满足其中两项时，本文才将其纳入路线比较或章节主线；若某材料仅以新闻稿、发布会口径、单次演示或未公开复现链路的技术验证出现，则只作为生态线索或趋势性观察引用，不直接据此给出成熟度与性能判断。

基于这一原则，本文将证据分为四层使用。第一层是系统论文与技术报告，用于界定执行、状态、压缩和评测等机制问题的研究脉络；第二层是官方文档、源码仓库与安装部署入口，用于确认对象的实际能力边界、组件关系与可进入性；第三层是标准、白皮书与生态报告，用于辅助刻画国产生态的分工变化与标准化趋势；第四层才是社区部署记录和兼容性经验，它们仅用于补充真实工程摩擦点，而不单独支撑横向性能结论。按这一分层，正文中的比较表与路线判断优先回答三类问题：一是复杂度主要被压在何处，二是哪些系统边界已经具有较稳定的公开证据，三是哪些判断仍应保留为条件性的工程观察。这样做的目的，不是把本文伪装为严格元分析，而是把综述中的描述性事实、比较性判断与趋势性推断尽量分层组织，使不同性质的材料不被直接并列。

\section{面向国产算力的大模型推理引擎生态}
\label{sec:domestic-ecosystem}

本章承接摘要提出的国产推理引擎分析框架，从产业生态与技术路线落地的维度，拆解国产算力场景下推理系统的完整技术栈构成。摘要将国产推理引擎的技术路线归纳为\textbf{开源主干后端适配、平台一体化方案与混合式架构}三类，本章的生态分层与技术实践恰好与三类路线形成一一映射：硬件厂商自研原生推理引擎对应平台一体化方案，主流开源推理框架的国产硬件适配对应开源主干后端适配路线，而中间层解耦与服务化平台则构成了混合式架构的核心载体。本章将从生态构成、代表实践、技术特征与演进趋势四个维度，系统梳理国产算力推理引擎的产业与技术现状，为后续的技术路径对比、核心技术拆解与架构演进分析提供生态层面的事实支撑\cite{caict2024whitepaper,zhang2024domestic}。

面向国产算力的大模型推理引擎生态并不是单一软件系统，而是由硬件厂商基础软件、开源推理框架适配层、算子与编译栈、以及上层服务化平台共同构成。根据公开仓库与官方文档，可以将相关工作大致归纳为三类：第一类是硬件厂商自研推理引擎或原生后端，例如华为 Ascend MindIE、寒武纪 MagicMind、摩尔线程 MT-Transformer、天数智芯 IxRT/IGIE 等；第二类是围绕 vLLM、SGLang、LMDeploy、FastDeploy 等主流大模型推理框架的国产硬件适配；第三类是 DLinfer、Triton-Ascend、Xinference、GPUStack 等中间层或服务化平台，用于降低上层框架、算子库、图编译器与硬件运行时之间的耦合。

三类生态实践对应了国产推理引擎研发的三条核心路径，其底层设计逻辑、技术取舍与目标场景存在本质差异，是本文后续进行技术选型对比的核心基础\cite{cui2025infer,sheng2024survey}：
\begin{enumerate}[leftmargin=*, label=(\arabic*)]
    \item \textbf{硬件厂商自研原生推理引擎（平台一体化路线）}：核心逻辑为“全栈协同优化”。向下深度绑定自研硬件架构与基础软件栈（如驱动、编译器、算子库），向上针对大模型推理的核心负载（如 Transformer 算子、KV Cache 管理、分布式通信）进行硬件原生定制优化。该路线能最大化释放硬件潜力，在 MoE 专家并行、低精度量化及内存层次优化等强硬件相关场景中优势显著；但其生态兼容性受限，通常仅支持特定硬件平台，开发者迁移成本较高。
    \item \textbf{主流开源推理框架的国产硬件适配（开源主干适配路线）}：核心逻辑为“生态兼容优先”。基于全球主流开源主干（如 vLLM、SGLang、LMDeploy 等），通过后端适配、算子移植与插件扩展，将成熟的开源能力迁移至国产硬件。该路线可最大化复用开源社区成果，降低研发成本并保持特性同步，支持多模型与多场景；但其硬件特性释放受限于框架抽象层，难以实现极致优化，且面临持续适配与版本碎片化的风险。
    \item \textbf{中间层解耦与服务化平台（混合式架构路线）}：核心逻辑为“解耦与标准化”。在推理框架与硬件后端之间构建统一抽象层，通过定义标准接口规范实现“一次开发、多硬件适配”。该路线有效解决了“N个框架$\times$M个硬件”的重复适配痛点，在异构算力集群与云原生规模化部署场景中具有显著优势；但其挑战在于统一抽象层的设计难度大，需在跨平台兼容性与硬件性能损耗之间寻求最优平衡。
\end{enumerate}

从公开生态看，vLLM 的硬件插件化正在成为国产算力适配的主要形态之一。vLLM-Ascend、vLLM-MLU、vLLM-MUSA、vLLM-MetaX、vLLM-Kunlun、vLLM-GCU 等项目均采用或参考 vLLM 的硬件插件机制，使厂商能够在尽量不侵入 vLLM 主干代码的情况下接入各自的 NPU、MLU、GPU、GCU 或 XPU 后端。这一路径的优势在于能够复用 vLLM 已有的连续批处理、KV Cache 管理、OpenAI 兼容服务接口、量化与投机解码等推理框架能力；挑战则在于底层算子、通信、图捕获、显存/片上内存管理、版本依赖和模型覆盖度仍需要各厂商持续适配。

vLLM 的硬件插件化架构经历了从\textbf{侵入式修改（In-Tree）}到\textbf{非侵入式插件（Out-of-Tree, OOT）}的关键演进，其核心技术突破在于构建了分层的硬件抽象接口：vLLM 在核心调度层、KV Cache 管理层、算子层分别定义了标准化的硬件无关抽象接口，厂商仅需实现对应接口的硬件后端即可完成接入，解决了早期适配中“一版本一修改、上游更新即失效”的维护难题。

当前，vLLM OOT 插件架构已成为国产算力适配的事实标准。国内主流硬件厂商均基于该架构完成了适配，同时行业内也在积极推进插件接口的标准化工作，如中国信通院牵头的《大模型推理引擎硬件适配接口规范》以及上海人工智能实验室 DeepLink 生态定义的统一算子接口\cite{kwon2024vllm,caict2025standard,shailab2024deeplink}。然而，当前插件化架构仍存在以下核心痛点：
\begin{itemize}[leftmargin=*, label=\textbullet]
    \item \textbf{标准化程度不足}：不同厂商的插件实现仍存在碎片化，难以实现跨硬件的插件复用。
    \item \textbf{硬件特性释放不充分}：MoE 分布式通信、片上内存管理等强硬件相关特性，仍需侵入式修改主干代码才能实现最优性能。
    \item \textbf{维护成本依然较高}：国产硬件适配代码大多维护在厂商独立仓库中，难以进入上游主干，面临长期同步成本。
\end{itemize}

华为 Ascend 生态是目前公开资料较为完整的代表。一方面，MindIE 作为面向 Ascend 的 AI 推理加速套件，覆盖大模型推理、服务化和 PyTorch 生态适配等能力；另一方面，vLLM-Ascend、SGLang-Kernel-NPU 和 Triton-Ascend 分别从推理框架插件、SGLang 算子库和 Triton 编译后端三个层面补充了开源生态。这说明国产算力推理引擎的竞争点已经从单一模型执行扩展到框架调度、KV Cache、MoE 通信、算子融合、编译优化和服务化接口等完整技术栈。

除单一芯片生态外，跨硬件适配层也在快速发展。上海人工智能实验室 DeepLink 生态中的 DLinfer 试图在上层 LLM 推理框架与下层厂商融合算子或图引擎之间定义统一接口，LMDeploy 则通过 PyTorch Engine 和 DLinfer 等路径支持多类国产硬件。百度飞桨 FastDeploy 以 PaddlePaddle 为基础，面向文心等大模型提供推理部署能力，并在官方文档中列出对昆仑芯、Ascend、海光、天数智芯、沐曦、燧原等硬件的适配说明。这类项目的共同目标是减少每个推理框架分别对接每个硬件后端的重复工程成本。

\begin{small}
\begin{longtable}{L{2.3cm}L{2.2cm}L{3.0cm}L{6.4cm}}
\caption{面向国产算力的大模型推理引擎与适配生态}
\label{tab:domestic-inference-engines} \\
\toprule
组织/单位 & 面向算力 & 代表项目/引擎 & 公开形态与技术关注点 \\
\midrule
\endfirsthead

\toprule
组织/单位 & 面向算力 & 代表项目/引擎 & 公开形态与技术关注点 \\
\midrule
\endhead

华为 Ascend 生态 & Ascend NPU，Atlas 推理/训练服务器 & MindIE，vLLM-Ascend，SGLang-Kernel-NPU，Triton-Ascend & 覆盖厂商推理套件、vLLM 硬件插件、SGLang NPU kernel 与 Triton 后端。技术关注点包括 CANN/torch-npu 适配、服务化推理、MoE kernel、注意力算子、编译后端和硬件插件化接口\cite{huawei2024mindie,huawei2024vllm_ascend,huawei2025sglang_npu,huawei2024triton_ascend}。 \\

SGLang 社区与 Ascend 适配 & Ascend NPU & SGLang-Kernel-NPU，DeepEP-Ascend & 面向 SGLang 的 Ascend NPU kernel 库，包含专家并行通信、attention、normalization、activation、LoRA 等推理关键算子，适合在综述中归入高性能 serving 框架的国产硬件 kernel 适配方向\cite{huawei2025sglang_npu}。 \\

上海人工智能实验室 DeepLink / OpenMMLab / InternLM 生态 & 多类国产加速器，包括 Ascend、Cambricon、MetaX 等 & DLinfer，LMDeploy & DLinfer 面向 LLM 推理框架 and 厂商算子/图引擎之间的解耦，LMDeploy 则提供 TurboMind 与 PyTorch Engine 两类推理引擎，并通过构建目标或适配层支持多种国产硬件\cite{shailab2024deeplink_repo,shailab2024dlinfer,shailab2025lmdeploy}。 \\

百度飞桨 / 昆仑芯生态 & Kunlun XPU，同时覆盖多种国产硬件 & FastDeploy，vLLM-Kunlun & FastDeploy 面向大模型和多模态模型的生产级部署，支持 PD 分离、统一 KV Cache 传输、OpenAI/vLLM 兼容接口、量化和投机解码等能力；vLLM-Kunlun 则以 vLLM 硬件插件形式支持昆仑 XPU\cite{baidu2025fastdeploy,kunlun2024vllm_kunlun}。 \\

寒武纪 & Cambricon MLU & vLLM-MLU，MagicMind，CNServing & vLLM-MLU 基于 vLLM 插件系统在 MLU 硬件上提供大模型推理服务能力；MagicMind 是寒武纪面向 MLU 的推理加速引擎，负责模型转换、图优化、代码生成和部署\cite{cambricon2024vllm_mlu,cambricon2024magicmind}。 \\

摩尔线程 & MUSA GPU & vLLM-MUSA，vLLM-MTT / MT-Transformer & vLLM-MUSA 强调与 vLLM 插件体系和 MUSA 软件栈集成；vLLM-MTT 则基于摩尔线程自研 MT-Transformer 后端，强调面向 Transformer 推理的底层算子融合和定制优化。二者体现了“通用兼容插件”和“高性能原生后端”两条路线\cite{moorethreads2024vllm_musa,moorethreads2024mt_transformer}。 \\

沐曦集成电路 & MetaX GPU，MACA 软件栈 & vLLM-MetaX，LMDeploy/DLinfer 适配 & vLLM-MetaX 通过 MACA 类 CUDA 后端接入 vLLM，目标是在沐曦 GPU 上获得接近 CUDA 生态的开发体验；DLinfer 也将 MetaX 作为其支持的硬件后端之一\cite{metax2024vllm_metax,shailab2024dlinfer}。 \\

燧原科技 & Enflame GCU，S60/L600 等 & vLLM-GCU，FastDeploy-Enflame & vLLM-GCU 面向燧原 GCU 适配 vLLM，支持量化、Fused MoE、Multi-LoRA、自动 Prefix Cache、Chunked Prefill、Speculative Decoding 等推理特性；FastDeploy 也提供燧原硬件部署说明\cite{enflame2024vllm_gcu,enflame2024fastdeploy_enflame,baidu2025fastdeploy}。 \\

天数智芯 / DeepSpark 生态 & Iluvatar CoreX、智铠/天垓系列 & IxRT，IGIE，DeepSparkInference，FastDeploy-Iluvatar & IxRT 和 IGIE 分别面向高性能推理运行时和基于 TVM 的通用推理引擎，DeepSparkInference 提供包括 ChatGLM、Baichuan、Qwen 等模型在内的推理示例；FastDeploy 也提供 Iluvatar CoreX 安装与部署路径\cite{iluvatar2024ixrt,iluvatar2024deepspark_infer}。 \\

海光信息 & Hygon DCU，K100-AI 等 & FastDeploy-Hygon DCU & 官方 FastDeploy 文档给出了在海光 DCU 上部署 ERNIE 系列模型的示例，说明该方向已有面向大模型推理的公开部署路径。但公开资料更多体现为部署适配和示例，成熟度仍需通过复现实验进一步判断\cite{hygon2024fastdeploy_hygon,baidu2025fastdeploy}。 \\

壁仞科技 & Biren SUPA / BR 系列加速器 & BIRENSUPA，EngineX-Biren vLLM 适配 & 壁仞公开开发者生态展示了 BIRENSUPA 软件平台；同时公开模型社区中存在基于 vLLM OOT 插件的壁仞适配项目。由于公开资料相对有限，综述中宜将其列为待持续跟踪的适配生态，而非直接给出性能结论\cite{biren2024supadoc,biren2024vllm_biren}。 \\

瀚博半导体 / Vastai 生态 & Vastai GPU，VACC/VUCA 软件栈 & Xinference VACC，Vastai 开发者平台 & 瀚博侧公开资料更多体现为开发者平台和模型生态；Xinference 等服务平台已出现对 Vastai GPU/VACC 的适配支持。该方向可归入国产硬件服务化部署平台生态，而不是单独的核心推理引擎\cite{vastai2024devdoc,xinference2025doc}。 \\

Xinference / GPUStack 等平台 & 多类异构 GPU/NPU & Xinference，GPUStack & 这类项目位于模型服务和集群管理层，通常不直接替代 vLLM、SGLang、LMDeploy 等推理引擎，而是对其进行编排、部署和多硬件管理。它们适合放在“服务化与运维平台”小节中讨论\cite{xinference2025doc,gpustack2024repo}。 \\

\end{longtable}
\end{small}

基于上述生态实践与技术特征分析，结合摘要提出的分析框架，本文从\textbf{生态兼容性、性能上限、复杂任务支持、运行时治理能力、长期维护成本}五个核心维度，对三类技术路线进行系统性对比，结果如表\ref{tab:comparison-matrix}所示。

\begin{table}[htbp]
\centering
\small
\caption{国产推理引擎三类技术路线多维度对比}
\label{tab:comparison-matrix}
\begin{tabularx}{\textwidth}{L{2.2cm}YYYY}
\toprule
对比维度 & 开源主干适配路线 & 平台一体化路线 & 混合式架构路线 \\
\midrule
生态兼容性 & 极高：完全兼容主流开源生态 & 较低：仅支持自研硬件 & 高：向上兼容主流框架 \\
单硬件性能上限 & 中：受限于框架抽象层 & 极高：全栈协同，原生优化 & 中高：平衡抽象与性能 \\
复杂任务支持 & 极高：社区特性同步快 & 中：依赖厂商研发周期 & 高：适配灵活度高 \\
运行时治理 & 中高：依赖框架原生能力 & 极高：全栈可控，定制深度 & 高：可基于统一底座构建 \\
长期维护成本 & 中高：需持续跟进上游 & 高：全栈自研，投入大 & 低：一次适配，多处复用 \\
核心适用场景 & 多硬件兼容、快速迭代 & 单硬件锁定、极致性能 & 异构集群、规模化部署 \\
\bottomrule
\end{tabularx}
\end{table}

该对比结果清晰地呈现了三类路线的技术取舍：不存在绝对最优的技术路线，仅存在与业务场景、硬件锁定程度、研发资源匹配的最优选型。这也呼应了本文核心观点：国产推理系统的核心竞争力不在于局部算子的极致性能，而在于技术路线与场景需求的匹配度，以及平台契约、运行时组织和交付闭环的协同能力\cite{caict2024whitepaper,zhang2024domestic}。

上述生态显示出几个趋势。首先，国产算力适配正在从修改上游推理框架源码转向硬件插件化和后端解耦。这种方式能够降低维护成本，也便于跟随 vLLM、SGLang 等主流推理框架的快速演进。其次，厂商原生推理引擎仍然重要，尤其是在算子融合、图优化、MoE 通信、低精度量化和特定硬件内存层次优化方面，原生后端通常更容易发挥硬件特性。再次，DLinfer、Triton-Ascend、FastDeploy 等中间层说明，国产算力生态正在从单点适配走向“统一接口 + 厂商实现”的模式。

在生态快速演进的同时，国产算力推理引擎仍面临以下三大核心挑战，这些挑战也构成了本文后续引入“一个底座、三条主路径”分析框架的出发点：
\begin{enumerate}[leftmargin=*, label=(\arabic*)]
    \item \textbf{生态碎片化与标准化的矛盾}：国产算力生态目前仍存在“一硬件一栈、一框架一适配”的碎片化现象。不同硬件平台的基础软件栈、算子接口与通信语义差异巨大，导致行业重复适配成本极高。这要求系统必须在架构底层构建统一的抽象底座，以吸收平台差异\cite{cesi2024standard}。
    \item \textbf{开源生态兼容与自主可控的矛盾}：主流推理框架体系仍由海外社区主导，国产适配多处于“跟随式”状态。如何在深度兼容开源生态的同时，构建自主可控的执行架构与状态治理能力，是掌握技术演进话语权的关键。
    \item \textbf{性能极致优化与工程交付闭环的矛盾}：现有优化多聚焦于单算子或单模型的峰值性能，而相对忽视生产环境下的工程交付闭环。生产级推理系统的核心竞争力在于执行、状态与成本三者的协同优化，而非实验室环境下的指标罗列\cite{qian2024hetero}。
\end{enumerate}

上述矛盾说明，单纯的“算子替换”已不足以支撑国产推理引擎的持续演进。为此，第 3 节将引入“一个底座、三条主路径”的分析骨架，从系统架构层面解构如何应对上述挑战。
\section{分析骨架与核心架构}

从系统视角看，大模型推理引擎是一套同时处理执行流、状态流与控制流的运行时。模型执行与算子层决定抽象模型如何映射到具体设备能力，缓存与调度层决定请求到达过程中资源如何被持续重写，接口与运维层则决定这些内部状态能否以可升级、可观测的方式进入真实生产。对国产平台而言，很多"局部性能"问题首先暴露在模块边界处：图模式回退是否需要被调度层感知，多模态前处理应停留在接口层还是进入状态系统，平台插件究竟只负责设备发现还是还要承担环境修复与后端选择。表\ref{tab:key-tech-points} 汇总了这一分析骨架下最核心的技术问题与评价维度。

\begin{table}[t]
\centering
\small
\caption{国产推理引擎关键技术点总览}
\label{tab:key-tech-points}
\begin{tabularx}{\linewidth}{L{2.5cm}YY}
\toprule
技术主路径 & 关键技术点 & 典型评价维度 \\
\midrule
执行与通信 & 统一执行 IR、动态图与图模式切换、Prefill/Decode 阶段解耦、推测解码协同、拓扑感知通信运行时、计算--通信重叠 & MFU、吞吐、TTFT、TPOT、跨节点扩展效率、回退稳定性 \\
状态治理 & Prefix/共享状态索引、PagedAttention 与块化缓存、多级存储驻留、状态生命周期管理、跨节点状态迁移、命中感知调度 & KV 命中率、迁移流量、恢复时延、显存/主存/NVMe 占用、长上下文稳定性 \\
压缩与成本协同 & 混合精度量化、KV 动态量化、稀疏--量化协同、低比特算子路径、质量约束下的单位 token 成本优化 & 精度保持、吞吐/时延变化、状态容量变化、带宽占用、单位 token 成本 \\
统一底座与交付 & 接口协议、指标口径、trace/log/metrics、回归验证、第三方测试证据链、环境治理与部署控制面 & 可复现性、可比较性、升级回归成本、交付成熟度 \\
\bottomrule
\end{tabularx}
\end{table}

\begin{figure}[!htb]
\centering
\resizebox{0.94\linewidth}{!}{\begin{tikzpicture}[
    node distance=5mm and 6mm,
    block/.style={draw, rounded corners=2pt, align=center, minimum height=9mm, inner sep=4pt, font=\small},
    group/.style={draw, rounded corners=3pt, inner sep=5pt},
    flow/.style={-{Latex[length=2.2mm]}, semithick}
]
\node[block, fill=gray!10, minimum width=2.9cm] (ingress) {业务请求\\文本/长上下文/多模态/工具调用};
\node[block, fill=blue!8, minimum width=2.7cm, right=of ingress] (service) {服务入口与接口层\\OpenAI 兼容/鉴权/监控};
\node[block, fill=blue!8, minimum width=2.7cm, right=of service] (scheduler) {调度与运行时层\\batch 组织/SLO/bucket/回退};
\node[block, fill=green!8, minimum width=2.4cm, below=of scheduler] (cache) {KV Cache 与状态层\\前缀复用/迁移/驱逐};
\node[block, fill=green!8, minimum width=2.4cm, left=of cache] (execute) {执行层\\prefill/decode/并行策略};
\node[block, fill=green!8, minimum width=2.4cm, right=of cache] (backend) {后端接口层\\图模式/算子/能力门控};
\node[group, fit=(execute)(cache)(backend), label={[font=\small]above:运行时核心}] (runtime) {};
\node[block, fill=orange!10, minimum width=2.6cm, right=18mm of backend] (compiler) {编译器与运行时\\CANN/MindSpore/插件栈};
\node[block, fill=orange!10, minimum width=2.3cm, below=of compiler] (device) {国产硬件平台\\NPU/GPU/驱动/通信};

\draw[flow] (ingress) -- (service);
\draw[flow] (service) -- (scheduler);
\draw[flow] (scheduler) -- (backend);
\draw[flow] (scheduler) -- (execute);
\draw[flow] (scheduler) -- (cache);
\draw[flow] (execute) -- (cache);
\draw[flow] (cache) -- (backend);
\draw[flow] (backend) -- (compiler);
\draw[flow] (compiler) -- (device);
\draw[flow] (device.west) |- (backend.east);
\draw[flow, dashed] (cache.west) |- (service.south);

\node[font=\scriptsize, align=left, below=2mm of ingress] {真实压力来自多阶段请求异质性};
\node[font=\scriptsize, align=left, below=2mm of compiler] {平台差异集中暴露在编译、算子与回退边界};
\end{tikzpicture}}
\caption{国产推理引擎系统设计示意图}
\label{fig:engine-architecture-overview}
\end{figure}

\subsection{从最小推理闭环到统一分析底座}

输入文本经 tokenizer 切分、embedding 转换后进入 Transformer block 序列。对推理服务而言，关键在于 prefill 与 decode 的行为差异：Prefill 对整段提示词并行建模并建立 KV Cache，随后 decode 每步只处理增量计算和缓存更新。首 token 时延主要受前处理与 prefill 主路径影响，稳定生成速度则更多受 decode 主路径与带宽约束影响。连续批处理、chunked prefill 和阶段解耦本质上都在利用前者更接近 compute-bound、后者更接近 memory-bound 这一差异\citep{orca2023,sarathi2023,zhong2024distserve}。

\begin{figure}[!htb]
\centering
\resizebox{0.97\linewidth}{!}{\begin{tikzpicture}[
    node distance=8mm and 10mm,
    block/.style={draw=black!75, line width=0.8pt, rounded corners=2.2pt, minimum width=25mm, minimum height=10mm, align=center, font=\small, inner sep=4pt},
    note/.style={font=\scriptsize, align=center, fill=white, inner sep=1pt, text=black!80},
    flow/.style={-{Latex[length=2mm]}, semithick, draw=black!75}
]
\node[block, fill=gray!8] (input) {输入文本\\tokenizer};
\node[block, fill=blue!9, right=of input] (embed) {embedding\\位置编码};
\node[block, fill=green!10, right=of embed] (prefill) {Prefill\\建立 KV Cache};
\node[block, fill=green!10, below=12mm of prefill] (decode) {Decode\\增量生成};
\node[block, fill=gray!8, right=of decode] (output) {输出 token\\流式返回};

\draw[flow] (input) -- (embed);
\draw[flow] (embed) -- (prefill);
\draw[flow] (prefill) -- (decode);
\draw[flow] (decode) -- (output);
\draw[flow, dashed] (prefill.south west) -- ++(-9mm,0) coordinate (ttft)
    -- (ttft |- decode.west) coordinate (ttft2) -- (decode.west);
\node[note, anchor=east, xshift=-2pt, text width=2.8cm] at ($(ttft)!0.5!(ttft2)$) {首 token 时延};

\node[note, below=8mm of decode, text width=5.2cm] {稳定生成速度主要受 decode 主路径、缓存组织和带宽利用影响。};
\end{tikzpicture}}
\caption{从输入到输出的最小推理闭环}
\label{fig:minimal-inference-loop}
\end{figure}

在此基础上，可观测底座的首要作用是为多平台、多版本和多阶段优化条件下的比较提供共同语义。若缺少统一指标口径（如 TTFT、TPOT、MFU 的对齐定义），不同路线的结论就无法在同一坐标系中比较；若缺少 trace 和 metrics 基础设施，调度策略的效果就难以在长时序中被持续验证。

\begin{figure}[!htb]
\centering
\resizebox{0.85\linewidth}{!}{\definecolor{obsBlue}{HTML}{DCEEFF}
\definecolor{obsGreen}{HTML}{DDF3E4}
\definecolor{obsOrange}{HTML}{FCE8D4}
\definecolor{obsGray}{HTML}{F3F4F6}

\begin{tikzpicture}[
    node distance=11mm and 15mm,
    core/.style={circle, draw=black!70, line width=0.9pt, minimum size=22mm, align=center, font=\small\bfseries, fill=obsGray},
    petal/.style={circle, draw=black!70, line width=0.8pt, minimum size=17mm, align=center, font=\small\bfseries},
    note/.style={font=\scriptsize\sffamily, text=black!72, align=center},
    link/.style={-{Latex[length=2.2mm,width=1.8mm]}, draw=black!70, line width=0.8pt},
    guide/.style={draw=black!40, dashed, line width=0.65pt}
]

\node[core] (core) {统一\\口径};

\node[petal, fill=obsBlue, above left=12mm and 20mm of core] (api) {接口\\语义};
\node[petal, fill=obsBlue, above right=12mm and 20mm of core] (metric) {指标\\口径};
\node[petal, fill=obsGreen, below left=12mm and 20mm of core] (state) {状态\\描述};
\node[petal, fill=obsOrange, below right=12mm and 20mm of core] (trace) {观测\\证据};

\node[note, above=3mm of api] {协议 / 版本};
\node[note, above=3mm of metric] {TTFT / TPOT};
\node[note, below=3mm of state] {命中 / 迁移};
\node[note, below=3mm of trace] {trace / log};

\draw[guide] (api) -- (core);
\draw[guide] (metric) -- (core);
\draw[guide] (state) -- (core);
\draw[guide] (trace) -- (core);

\draw[link] (api.east) to[bend left=16] (metric.west);
\draw[link] (metric.south) to[bend left=16] (trace.north);
\draw[link] (trace.west) to[bend left=16] (state.east);
\draw[link] (state.north) to[bend left=16] (api.south);

\node[draw=black!60, rounded corners=2.2pt, fill=white, inner sep=4pt, align=center, font=\scriptsize\sffamily, below=21mm of core, minimum width=58mm] (outcome) {路线比较 \quad $\rightarrow$ \quad 回归验证 \quad $\rightarrow$ \quad 交付诊断};
\draw[link] (core) -- (outcome);

\node[note, right=12mm of metric, align=left] {统一后才能\\跨平台比较};
\node[note, left=12mm of state, align=right] {把指标、状态与\\日志放进同一证据链};

\end{tikzpicture}}
\caption{可观测性与能力契约层级示意图}
\label{fig:observability-contract-stack}
\end{figure}

\subsection{主路径一：执行与通信}

执行路径决定推理引擎的吞吐上限与时延下限。核心问题包括：何时使用图模式、何时回退到 eager 执行；如何在多设备间组织 tensor parallel、pipeline parallel 及 MoE expert parallel 的通信拓扑；以及如何在推测解码、chunked prefill 和阶段解耦中平衡计算与通信重叠。

对国产平台，执行路径还包含接口稳定性与维护成本。如果后端为获得局部加速在共享路径中嵌入大量平台条件分支，一旦上游框架修改模型执行接口、缓存布局或采样流程，本地后端就会承受高维护负担。FlashAttention-3 直接利用 H100 上 Tensor Core 异步执行与 FP8 支持，FlexAttention 则在保留 fused kernel 性能的同时恢复 attention 变体的可编程性\citep{shah2024flashattention3,dong2025flexattention}。这类工作说明，高性能后端的难点往往不在"有没有 kernel"，而在"如何把硬件代际能力组织成可持续复用的编译与运行时接口"。

因此，面向国产算力的设计应优先使用平台注册表、设备抽象和能力门控，把"平台差异"封装成可替换模块。算子和图编译层更需要沉淀的是一套后端能力契约：哪些算子变体可被图模式安全覆盖，哪些路径必须保留 eager 回退，哪些低比特逻辑需要运行时显式感知。对国产平台而言，很多问题并不是"后端不可用"，而是"某种能力只在某些约束下可用"。若系统无法表达这种带条件的能力边界，调度器就只能以越来越多的探测和分支维持稳定。

\begin{figure}[!htb]
\centering
\resizebox{0.88\linewidth}{!}{\begin{tikzpicture}[
    node distance=10mm and 12mm,
    block/.style={draw=black!75, line width=0.8pt, rounded corners=2.2pt, minimum width=30mm, minimum height=10.5mm, align=center, font=\small, inner sep=4pt},
    flow/.style={-{Latex[length=2mm]}, semithick, draw=black!75},
    note/.style={font=\scriptsize, align=center, fill=white, inner sep=1pt, text=black!80}
]
\node[block, fill=blue!9] (scheduler) {执行编排\\batch / 阶段切换};
\node[block, fill=green!10, right=of scheduler] (runner) {model runner\\worker / 并行策略};
\node[block, fill=orange!12, below=12mm of runner] (backend) {后端能力\\图模式 / kernel / 回退};
\node[block, fill=gray!8, right=of backend] (device) {设备与通信\\NPU / GPU / 驱动};

\draw[flow] (scheduler) -- (runner);
\draw[flow] (runner) -- (backend);
\draw[flow] (backend) -- (device);
\draw[flow, dashed] (device.north) |- (runner.east);

\node[note, below=8mm of backend, text width=6.2cm] {平台差异更适合被封装在后端能力边界内，而不是持续回流到共享调度主链。};
\end{tikzpicture}}
\caption{执行后端能力边界与回退路径}
\label{fig:execution-backend-boundary}
\end{figure}

\subsection{主路径二：状态治理与 KV Cache 组织}

KV Cache 及其元数据系统是状态治理的核心载体，直接影响并发能力、会话恢复成本与长上下文性能。PagedAttention 通过块化管理显存，降低了连续大块分配的压力，也为前缀复用、分层缓存与换入换出提供了实现基础\citep{kwon2023efficient}。对国产硬件，状态系统面临两个额外约束：不同后端对内存分配粒度与 host-device 同步开销的敏感度不同；当系统需要支持 Prefix Cache、Chunked Prefill 和跨设备迁移时，缓存元数据管理复杂度显著提升。

围绕注意力高效实现与缓存压缩已形成清晰技术主线。FlashAttention 系列从算子层逐步降低注意力计算的内存访问代价\citep{dao2022flashattention,dao2024flashattention2,shah2024flashattention3}；MInference、DuoAttention、vAttention 与 LServe 分别从动态稀疏注意力、检索头/流式头分离、虚拟内存管理和统一稀疏注意力执行改写了长上下文推理的系统组织\citep{jin2024minference,xiao2024duoattention,kwon2024vattention,yang2025lserve}；H2O、KIVI、SnapKV 与 CacheGen 则从重点击中、低比特量化、语义压缩和缓存流式传输角度缓解了 KV Cache 的显存与带宽压力\citep{zhang2023h2o,liu2024kivi,li2024snapkv,zhang2024cachegen}。

从系统演进角度看，KV Cache 正从执行层附属模块演变为独立基础设施。随着长上下文、多轮会话和多 Agent 工作流共存，缓存不仅要回答"数据放在哪"，还要回答"状态如何被计价、回收、复用和转移"。国产推理引擎在 KV Cache 设计上更适合采用"统一抽象 + 平台特化实现"的方式：统一暴露块管理、引用计数、前缀复用和驱逐策略接口，而将具体数据排布与搬运实现留给后端适配层。

\begin{figure}[!htb]
\centering
\resizebox{0.58\linewidth}{!}{\begin{tikzpicture}[
    node distance=12mm and 16mm,
    block/.style={draw=black!75, line width=0.8pt, rounded corners=2.2pt, minimum width=28mm, minimum height=10mm, align=center, font=\small, inner sep=4pt},
    flow/.style={-{Latex[length=2mm]}, semithick, draw=black!75}
]
\node[block, fill=green!10] (cache) {KV Cache\\块化 / 前缀复用};
\node[block, fill=orange!12, right=of cache] (compress) {压缩路径\\量化 / 稀疏 / 低比特};
\node[block, fill=blue!9, below=of compress] (cost) {成本结果\\显存 / 带宽 / 单位 token};
\node[block, fill=gray!8, left=of cost] (quality) {服务质量\\吞吐 / 时延 / 精度};

\draw[flow] (cache) -- (compress);
\draw[flow] (compress) -- (cost);
\draw[flow] (cost) -- (quality);
\draw[flow] (quality) -- (cache);
\end{tikzpicture}}
\caption{KV Cache、压缩与成本反馈闭环示意图}
\label{fig:kv-cost-feedback-loop}
\end{figure}

\subsection{主路径三：压缩、成本与服务质量协同}

在国产平台语境下，压缩加速不应被看作部署后期附加的"省显存技巧"，而应被视为与执行路径和状态治理并列的核心主路径。混合精度量化、KV 动态量化和稀疏--量化协同确实能直接降低显存占用与单位 token 成本，但其实际收益高度依赖后端是否存在稳定低比特路径、状态层是否允许容量变化、调度层是否能吸收 batch 形态变化。也就是说，压缩配置的效果从来不由模型精度单独决定，而由"压缩参数、执行效率、状态占用、服务质量"四者共同决定。

如图\ref{fig:kv-cost-feedback-loop} 所示，压缩路径的工程价值需要回到 KV 组织、服务质量和单位 token 成本的联动关系中才能被正确评价。QServe 通过把权重量化、KV 压缩与 serving 系统协同组织，展示了低比特部署必须在 kernel、调度和缓存之间联合设计\citep{wang2024qserve}；KIVI、SnapKV 则进一步表明，长上下文场景下的 KV 压缩只有和缓存布局与热度判断结合，才可能不显著破坏服务质量\citep{liu2024kivi,li2024snapkv}。国产推理引擎若要把"低成本"转化为系统能力，就必须让压缩配置、运行时观测与成本核算使用统一接口。

\subsection{服务层收束：调度、接口与交付闭环}

\begingroup
\sloppy
服务层需要在吞吐与时延之间平衡。典型机制包括连续批处理、请求优先级控制、SLA 感知调度和流式返回。对国产硬件，还需考虑编译图缓存、静态形状约束与后端初始化成本——调度不仅要优化"请求何时执行"，还要优化"请求如何映射到平台最擅长的执行形态"。XGrammar 说明结构化生成已需要与推理引擎共同设计\citep{dong2024xgrammar}；AI Metropolis 进一步表明，多 Agent 负载下依赖感知的乱序执行是提高并行度的关键机制\citep{xie2025aimetropolis}。

在调度策略方面，DistServe 和 Splitwise 将 Prefill 与 Decode 的资源需求进一步拆解\citep{zhong2024distserve,splitwise2024}；Llumnix 通过跨实例动态迁移请求缓解负载失衡，SOLA 则把调度目标显式对齐到 TTFT 和 TPOT 等服务级目标\citep{sun2024llumnix,hong2025sola}。TensorRT-LLM 已把多模态支持做成相对完整的运行时能力集合\citep{tensorrtllmmultimodal2026}；OmniServe 尝试把低比特量化与长上下文统一稀疏注意力纳入同一 serving 框架\citep{omniserve2025}。这些工作虽不直接针对国产算力，但为理解调度协同、量化部署和工程化取舍提供了对照。

\endgroup

状态治理推进到生产环境后，最终落到交付闭环。推理引擎的对外层不仅包括 OpenAI 兼容接口，也包括监控、日志、链路追踪和故障恢复。这一层的重要性在国产推理系统中经常被低估，但在政企部署中，系统往往首先因兼容性、可观测性或升级复杂度而被淘汰。无论是 Transformers 的模型接口、Ray Serve 的服务编排，还是 Triton Inference Server 的部署体系，都在事实上塑造了推理引擎需要兼容的工程边界\citep{transformers,rayserve,tritonserver}。

\begin{figure}[!htb]
\centering
\resizebox{0.92\linewidth}{!}{\begin{tikzpicture}[
    node distance=16mm and 18mm,
    >=Stealth,
    mainblock/.style={
        draw=black!85, 
        line width=1pt, 
        rounded corners=3pt, 
        minimum width=38mm, 
        minimum height=14mm, 
        align=center, 
        fill=gray!8, 
        inner sep=4pt
    },
    opsblock/.style={
        mainblock,
        fill=white,
        draw=black!70,
        line width=0.75pt,
        dashed
    },
    flow/.style={->, line width=1pt, draw=black!90},
    feedback/.style={->, line width=0.75pt, dashed, draw=black!70},
    boundary/.style={
        draw=black!40, 
        line width=0.5pt, 
        dashed, 
        rounded corners=4pt, 
        fill=gray!2, 
        inner sep=18pt
    },
    edgelabel/.style={
        font=\scriptsize\sffamily, 
        text=black!75, 
        fill=gray!2, 
        inner sep=2pt
    }
]

\node[mainblock] (sched) {\textbf{调度与流量整形}\\[1.5mm]\footnotesize 令牌桶 / SLO};
\node[mainblock, left=of sched] (request) {\textbf{请求进入}\\[1.5mm]\footnotesize 服务接口};
\node[mainblock, right=of sched] (exec) {\textbf{执行与返回}\\[1.5mm]\footnotesize 流式输出};

\node[opsblock, below=18mm of sched] (ops) {\textbf{观测与运维}\\[1.5mm]\footnotesize 日志采集 $\vert$ 性能监控 $\vert$ 异常回退};

\path (exec.east) ++(28mm, 0) coordinate (right_padding);

\begin{scope}[on background layer]
    \node[boundary, fit=(request) (exec) (ops) (sched) (right_padding)] (system_boundary) {};
\end{scope}

\node[draw=black!40, line width=0.5pt, fill=white, rounded corners=2pt, inner sep=4pt, 
      anchor=south east, xshift=-4mm, yshift=4mm] at (system_boundary.south east) {
    \begin{tabular}{@{}cl@{}}
    \tikz[baseline=-0.5ex]\draw[flow] (0,0) -- (0.4,0); & \scriptsize 主请求流 \\
    \tikz[baseline=-0.5ex]\draw[feedback] (0,0) -- (0.4,0); & \scriptsize 运维反馈流
    \end{tabular}
};

\draw[flow] (request) -- (sched);
\draw[flow] (sched) -- (exec);
\draw[feedback] (exec.south) |- node[pos=0.75, edgelabel, anchor=south] {监控数据采集} (ops.east);
\draw[feedback] (ops.west) -| node[pos=0.25, edgelabel, anchor=south] {异常回退 / 配置调整} (request.south);

\end{tikzpicture}}
\caption{调度、服务与运维闭环示意图}
\label{fig:service-ops-closure}
\end{figure}

综合来看，国产推理引擎的架构难点集中在模块交界处：执行层与图编译层决定可运行边界，缓存与调度层决定长上下文和并发表现，对外接口与运维层决定系统能否形成生产闭环。后文在讨论典型技术路线时，重点关注的正是不同系统如何在这些边界上做取舍。

\section{典型技术路线与复杂度分配}

如果回归到工程师面对的真实决策场景，路线选择的核心问题会变得非常具体：当时间、人力和平台约束都有限时，团队究竟应该把有限的工程力量优先投入在哪一层的优化上？国产推理引擎当前的三类路线，正好对应了三种不同的答案。

第一种选择是把力量压在执行主路径上：以最快的速度让一个成熟的开源框架在国产后端上跑起来，复用上游社区在调度、缓存和接口方面的持续演进。这条路线的回报周期最短，但风险也最明确——上游每一次版本升级都可能在插件边界制造新的摩擦。第二种选择是把力量压在平台栈的整体收拢上：由厂商将推理引擎、编译器、运行时和部署工具打包为一套端到端方案，追求确定性交付和深度调优。这条路线的交付质量最高，但前提是团队愿意接受平台定义的边界和节奏。第三种选择则试图在两者之间找平衡：在保持上层接口兼容的前提下，通过插件隔离和外部治理工具把国产平台的适配劳动从核心执行链中拆开，既不完全依赖上游，也不被单一平台绑定。

表\ref{tab:community-route-selection} 将这三类路线的适用前提、典型摩擦和系统收益做了对比。但表格只能呈现静态快照，实际决策还需要理解每条路线内部的张力演化——这正是以下各小节试图展开的内容。

\subsection{路线一：把复杂度压在执行主链的开源后端适配}

以 vLLM 与 SGLang 为代表的开源推理框架，正通过平台接口、后端插件和定制算子适配国产硬件\citep{vllm,sglang,sglangpaper,vllmascend}。该路线的优势在于持续复用上游社区在调度、缓存、结构化输出与多模态支持方面的能力演进，并借助已有模型生态快速覆盖主流开源模型。典型实现方式包括：在平台层增加设备选择与能力注册；在 worker 层注入后端特化逻辑；在注意力、归一化等热点路径增加定制算子；以及在调度与缓存层引入面向本土平台的约束门控。不足在于，若抽象边界设计不当，后端 patch 会随上游快速演进不断累积技术债。

在国内部署实践中，这一路线的真实工作量不止于补齐 kernel，还包括模型接入、量化转换、OpenAI 兼容协议、监控埋点和部署脚本等外围工程。后端适配是否成功，往往取决于模型支持矩阵更新速度、异常回退路径是否可观测，以及版本升级时能否继续跟上上游抽象变化。vLLM-Ascend 更能体现对上游执行边界的依赖，LMDeploy 更强调部署工具链与模型接入效率，而 vLLM-HUST 则在上层保持 vLLM 兼容的同时，通过平台插件与外部运行时治理工具把国产平台接入从核心执行链中拆开\citep{vllmascend,lmdeploy,vllmhust}。这种差异说明，开源主干后端适配内部也在分化为"主链后端实现优先"和"模块边界经营优先"两种形态。

\subsection{路线二：把复杂度压在交付链路的一体化方案}

以 MindIE 为代表的路线，把推理引擎、编译器、运行时与部署工具链整体收拢到同一交付体系内\citep{mindie,huawei_mindie_2026}。其核心理念是系统应从编译到上线在相同约束空间中统一治理。这一路线在确定性交付和稳定迭代方面具有优势：模型覆盖由平台统一规划，算子融合深度不受通用 API 约束，服务配置可与平台基础设施联动。不足在于生态兼容与上游响应速度相对较慢，模型类型扩展受制于平台研发节奏，且因公开资料有限难以从外部独立验证能力边界。

\subsection{路线三：自研引擎与混合式架构}

Chitu（赤兔）与 xLLM 代表了另一种路径：围绕自研状态机、缓存治理和部署闭环组织系统能力\citep{chitu_github,xllm_docs,xllm_tech_report}。前者面向多元算力组织弹性推理部署，后者以国产芯片为中心组织推理内核并以前端部署入口补齐一键拉起与 OpenAI 兼容能力。这类路线的优势在于系统边界自主可控、状态治理可内生演进；代价则是生态吸收较慢、研发投入较高。

\begin{table*}[t]
\centering
\footnotesize
\setlength{\tabcolsep}{4pt}
\caption{国产推理引擎路线选择的前提、摩擦与收益}
\label{tab:community-route-selection}
\begin{tabularx}{\textwidth}{L{2.2cm}L{3.5cm}L{3.5cm}L{3.0cm}Y}
\toprule
路线 & 适用前提 & 典型摩擦 & 核心系统收益 & 适合的团队画像 \\
\midrule
vLLM-Ascend 类开源后端适配 & 愿意跟进上游主干演进，优先保障模型覆盖与生态兼容 & 版本矩阵对齐、patch 累积、平台特例向共享主链渗透\citep{vllmascend} & 快速吸收上游能力，模型覆盖广 & 重视生态距离和快速跟进的团队 \\
MindIE 类一体化平台 & 面向确定性交付，接受平台约束换取稳定性 & 生态兼容慢、模型类型受平台节奏限制\citep{mindie} & 交付确定性高，模板治理统一 & 政企落地与稳定交付优先的团队 \\
SGLang 路线 & 在复杂控制流和结构化输出场景下重视运行时组织优势 & 需为执行模型与调度抽象投入理解成本\citep{sglang,sglangpaper} & 复杂工作负载调度与编排 & 复杂工作负载优先的团队 \\
vLLM-HUST 类混合式路线 & 不愿放弃上游接口兼容，同时希望逐步沉淀本地治理能力 & 需持续经营模块边界，避免继承上游复杂度又堆积本地补丁\citep{vllmhust} & 接口兼容、平台特化和运行时治理之间的折中 & 既要跟进上游又要建设本地化治理的团队 \\
\bottomrule
\end{tabularx}
\end{table*}

\subsection{多模态与复杂负载对路线分化的加速}

如果说前面的讨论还主要围绕文本推理展开，那么多模态和复杂负载的到来则在实质上加速了路线的分化。问题的本质不在于"支持更多数据类型"，而在于当请求从单一文本扩展为包含视觉编码、语音前端、结构化约束验证和工具调用回注等多阶段链路时，三条主路径上的矛盾会同时加剧。

具体而言，多模态推理在三个层次上重写了引擎的假设。首先是 prefill 主路径：多图文档输入导致视觉 token 膨胀和形状剧烈波动，原本为文本 batch 优化的 bucket 划分和编译策略会频繁失效。其次是阶段关系本身：视频理解要求 encode 和 decode 之间插入帧间压缩步骤，文本推理时代"一条请求一个 decode 循环"的模型被打破。第三是会话状态机：语音交互引入流式 chunk 和实时打断，使引擎必须同时管理文本、音频和视觉多种状态对象，而非只维护一个 token 序列。

这些变化对三条路线的影响并不对称。对开源后端适配路线，冲击集中在执行主路径上——每次新模态出现都需要评估后端算子的覆盖度、编译路径的完整性和异常回退的可靠性，而这些正是此类路线最脆弱的环节。对一体化方案，挑战则在生态响应速度——当模型形态以月为单位快速变化时，平台统一规划的模式可能难以同步跟进。对混合式路线，多模态反而为其提供了合法性的论据：既然没有任何单一路线能覆盖所有模态的最优路径，那么"经营模块边界、在关键路径上允许平台特化"就变得比"追求全栈统一"更务实。这一判断也解释了为什么像 vLLM-HUST 这类路线近年受到关注——它不是某种技术偏好的产物，而是多模态时代对系统架构提出的实际要求。

\subsection{路线比较与选择框架}

\begin{table}[t]
\centering
\scriptsize
\setlength{\tabcolsep}{4pt}
\caption{国内部署路线案例比较}
\begin{tabularx}{\linewidth}{L{2.0cm}L{1.7cm}L{2.7cm}Y}
\toprule
代表对象 & 路线类型 & 工程落点 & 部署目标 \\
\midrule
vLLM-Ascend\citep{vllmascend} & 开源主干后端适配 & 围绕上游接口扩展后端与定制算子 & 生态兼容与模型覆盖优先 \\
vLLM-MLU\citep{vllmmlu} & 开源主干后端适配 & 依托 vLLM plugin system 接入寒武纪 MLU & 继承 vLLM 接口与调度语义 \\
SGLang Ascend\citep{sglang_ascend_npu} & 主线后端支持 & 在 SGLang 主线中为 Ascend NPU 提供 attention backend & 国产 NPU 被主流框架吸纳 \\
vLLM-HUST\citep{vllmhust} & 混合式架构 & 保持上层接口，外挂平台插件与运行时护栏 & 上游兼容、国产平台接入与运维闭环 \\
MindIE\citep{mindie} & 一体化部署 & 与本地编译器、运行时深度协同 & 平台统一交付优先 \\
Chitu\citep{chitu_github} & 独立自研引擎 & 面向多元算力组织弹性部署 & 多硬件兼容与生产可用性 \\
xLLM\citep{xllm_docs} & 自研框架 & 围绕国产芯片组织推理内核与部署入口 & 同步建设推理内核与部署闭环 \\
LMDeploy\citep{lmdeploy} & 工具链强化 & 围绕模型支持、推理工具链持续扩展 & 高性能部署与工程效率 \\
FastDeploy\citep{fastdeploy} & 工具链强化 & 以高性能工具包衔接模型导出与服务封装 & 围绕飞桨生态组织 LLM 部署 \\
\bottomrule
\end{tabularx}
\end{table}

\begin{figure}[!htb]
\centering
\resizebox{0.94\linewidth}{!}{\begin{tikzpicture}[
    node distance=5mm and 10mm,
    block/.style={draw, rounded corners=2pt, minimum width=31mm, minimum height=9mm, align=center, font=\small, inner sep=4pt},
    note/.style={draw, rounded corners=2pt, minimum width=31mm, align=center, font=\scriptsize, inner sep=3pt},
    flow/.style={-{Latex[length=2.2mm]}, semithick}
]
\node[block, fill=gray!10, minimum width=42mm] (goal) {国产推理引擎路线选择};
\node[font=\scriptsize, align=center, above=1mm of goal] {核心张力：生态兼容性、平台特化深度、交付稳定性};

\node[block, fill=blue!10, below left=12mm and 31mm of goal] (upstream) {开源主干后端适配\\vLLM / SGLang / vLLM-Ascend};
\node[block, fill=orange!12, below=12mm of goal] (hybrid) {混合式架构\\vLLM-HUST\\上层兼容 + 底层热点特化};
\node[block, fill=green!10, below right=12mm and 31mm of goal] (integrated) {一体化方案\\MindIE / 编译器 / 运行时协同};

\node[note, fill=blue!5, below=of upstream] (u1) {优势：生态吸收快\\模型跟进快};
\node[note, fill=red!5, below=of u1] (u2) {代价：patch 累积\\抽象边界易失稳};

\node[note, fill=orange!6, below=of hybrid] (h1) {优势：兼顾兼容性\\与平台优化};
\node[note, fill=red!5, below=of h1] (h2) {代价：体系复杂\\组织门槛高};

\node[note, fill=green!5, below=of integrated] (i1) {优势：交付稳定\\平台约束强};
\node[note, fill=red!5, below=of i1] (i2) {代价：生态兼容慢\\上游响应较慢};

\draw[flow] (goal) -- (upstream);
\draw[flow] (goal) -- (hybrid);
\draw[flow] (goal) -- (integrated);
\draw[flow] (upstream) -- (u1);
\draw[flow] (u1) -- (u2);
\draw[flow] (hybrid) -- (h1);
\draw[flow] (h1) -- (h2);
\draw[flow] (integrated) -- (i1);
\draw[flow] (i1) -- (i2);
\end{tikzpicture}}
\caption{国产推理引擎技术路线分化示意图}
\label{fig:route-divergence}
\end{figure}

将上述分析放在一起，可以提炼出一个更加底层的判断框架。路线比较的本质不是判断"A 比 B 好"，而是判断"在给定的约束下，哪一个架构维度的优先级应该排在最前面"。如果核心约束是模型覆盖速度和生态兼容性，那么优先保证执行主路径畅通、将其他维度暂交给社区和厂商分别解决的路线更务实。如果核心约束是端到端交付质量——在政企场景中极为常见——那么将执行、编译和运维一体化治理的路线更有优势，尽管它需要接受更慢的模型跟进节奏。如果核心约束是"既要跟进上游又要建设本地化能力"——这正是许多国内团队的真实处境——那么混合式路线提供了一种可行的折中，但其持续成立的前提是能够在快速变化的上游接口与相对稳定的本地治理需求之间维持清晰的模块边界。

除路线本身的选择外，还有一个容易被忽略的维度正在成为新的分化点：多租户与复杂状态空间下的运行时组织能力。当同一个引擎实例需要同时服务数十个 LoRA 适配器、管理多个模型版本的 KV 共享内存池、并在不同优先级的多 Agent 工作流之间动态分配计算资源时，"调度的粒度"就不再是 batch 尺寸调优，而是运行时是否具备对状态对象进行统一建模和动态编排的能力。这组问题超出了单一框架或平台的范围，直接指向下一节要讨论的核心——当这些架构选择在国产平台的实际部署中相互碰撞时，哪些结构性挑战会首先暴露出来。

\section{关键挑战：主路径如何相互放大}

国产推理引擎面临的难点并非单点性能不足，而是多个系统目标同时受限：既要持续吸收上游能力演进，又要适应本土硬件与基础软件栈的边界，还要在真实业务负载下维持稳定交付。这些挑战之所以顽固，是因为它们会相互放大——平台差异抬高抽象成本，抽象不稳又削弱复杂任务支持；复杂任务若采用侵入式补丁实现，又加速本地分叉累积；而如果评测仍主要停留在理想化 benchmark，团队就难以及时识别这些结构性问题。表\ref{tab:challenge-map} 将这种放大关系压缩为总览。

\begin{table}[t]
\centering
\small
\caption{国产推理引擎关键挑战比较}
\label{tab:challenge-map}
\begin{tabularx}{0.88\linewidth}{L{2.2cm}L{4.2cm}Y}
\toprule
挑战来源 & 主要冲击层 & 典型后果 \\
\midrule
硬件异构与后端碎片化 & 执行层、算子接口、通信路径 & 共享路径分叉、正确性边界不一致、维护成本持续上升 \\
复杂任务与多模态负载 & 调度、缓存、图模式与前处理链路 & 阶段争用加剧、时延波动放大、端到端不稳定 \\
压缩收益与系统成本错配 & 量化路径、KV 状态、调度与成本核算 & 低比特收益不稳定、单位 token 成本难形成闭环 \\
上游快速演进与本地分叉 & 平台抽象、插件接口与共享模块 & patch 累积、合并困难、对新模型和新特性的跟进滞后 \\
评测体系与真实业务脱节 & benchmark 设计、运维指标与资源评估 & 优化方向失真、离线结论与生产表现偏离 \\
\bottomrule
\end{tabularx}
\end{table}

\begin{table*}[t]
\centering
\footnotesize
\setlength{\tabcolsep}{4pt}
\caption{不同国产推理路线与关键挑战的映射关系}
\label{tab:route-challenge-mapping}
\begin{tabularx}{\textwidth}{L{2.1cm}L{2.8cm}L{2.5cm}L{2.8cm}L{2.8cm}Y}
\toprule
代表路线 & 版本矩阵/环境门槛 & 网络与分布式风险 & 缓存/调度压力 & 运维与回退焦点 & 更需要优先建设的系统能力 \\
\midrule
vLLM-Ascend 类开源后端适配 & 高。上游主干、插件分支、CANN 与 torch\_npu 对齐要求强\citep{vllmascend} & 中到高。PD 分离与多机连通性常在扩容时暴露 & 中。继承上游调度能力，但平台约束可能导致命中率下降 & 版本锁定与环境预检 & 持续回归基线与插件边界稳定性 \\
MindIE 类一体化平台 & 中。平台内控但对外兼容性需额外验证\citep{mindie} & 低到中。平台内一体治理 & 中。内部调度但缺乏外部可观测证据 & 模板敏感性与配置可控性 & 服务模板的可诊断性与生态兼容接口 \\
Chitu / xLLM 类自研路线 & 中。自主控制但生态距离远\citep{chitu_github,xllm_docs} & 取决于自研通信层成熟度 & 高。状态治理需自建 & 全栈自主但生态覆盖有限 & 模型覆盖速度与社区生态对接 \\
vLLM-HUST 类混合式路线 & 中到高。继承上游版本压力，同时维护本地插件层\citep{vllmhust} & 中。继承上游分布式能力，补充平台护栏 & 中。上游调度加本地门控 & 模块边界维护与 merge-safe 演进 & 持续经营模块边界防止双重分叉 \\
\bottomrule
\end{tabularx}
\end{table*}

表\ref{tab:route-challenge-mapping} 进一步将这种放大关系映射到具体路线上，说明不同路线并非面对同一组问题再做不同优化，而是会把“版本矩阵、网络控制面、缓存/调度、运维回退”这些挑战优先暴露在不同边界上。路线选择本身，也就在决定团队首先要把系统能力压在哪一层。

公开社区部署记录与工程经验反复印证了这一判断：CANN/torch\_npu 版本矩阵对齐困难（vLLM-Ascend 0.7.x 到 0.8.x 升级期间多次出现 HCCL 通信与算子兼容性问题）、PD 分离与双机组网的稳定性、MindIE 服务化配置可控性，以及 LMDeploy 与 vLLM 在国产 GPU 环境中的适配取舍——这些反复出现的摩擦点，最终都能回收到表\ref{tab:challenge-map} 列出的几类结构性挑战上。真实摩擦往往先出现在环境治理和平台工具链层面，而不是单点跑分上。下面沿前文三条主路径，逐一讨论这些挑战如何沿系统骨架扩散。

\begin{figure}[!htb]
\centering
\resizebox{0.80\linewidth}{!}{\definecolor{cBlue}{HTML}{E6F0FA}
\definecolor{cGreen}{HTML}{E9F5EE}
\definecolor{cOrange}{HTML}{FCEFE3}
\definecolor{cRed}{HTML}{F7E6E6}
\definecolor{cGray}{HTML}{F2F3F5}
\begin{tikzpicture}[
    node distance=11mm and 14mm,
    scale=1.10,
    transform shape,
    challenge/.style={draw=black!70, line width=0.8pt, rounded corners=2.2pt, minimum width=3.0cm, minimum height=10mm, align=center, font=\small, inner sep=4pt},
    flow/.style={-{Latex[length=2.5mm]}, line width=0.85pt, draw=black!70},
    amplify/.style={-{Latex[length=2.5mm]}, line width=0.9pt, draw=red!65, dashed},
    note/.style={font=\footnotesize, align=left, text=black!70}
]

% Center: unified base
\node[challenge, fill=cGray, minimum width=3.8cm, minimum height=10mm] (base) {\textbf{统一底座}\\上游演进与评测体系};

% Four challenges around
\node[challenge, fill=cBlue, above left=13mm and 4mm of base] (c1) {\textbf{平台异构扩散}\\执行路径持续分化};
\node[challenge, fill=cGreen, above right=13mm and 4mm of base] (c2) {\textbf{复杂任务冲击}\\调度/缓存/状态失控};
\node[challenge, fill=cOrange, below left=13mm and 4mm of base] (c3) {\textbf{压缩收益失调}\\离线收益难以兑现};
\node[challenge, fill=cRed, below right=13mm and 4mm of base] (c4) {\textbf{评测体系脱节}\\优化方向与实际背离};

% Flows from challenges to base (challenges feed into base)
\draw[flow] (c1) -- (base);
\draw[flow] (c2) -- (base);
\draw[flow] (c3) -- (base);
\draw[flow] (c4) -- (base);

% Amplification arrows between challenges
\draw[amplify, bend left=15] (c1) to node[above, font=\scriptsize, text=red!70] {平台差异抬高抽象成本} (c2);
\draw[amplify, bend left=15] (c2) to node[right, font=\scriptsize, text=red!70] {侵入补丁加速分叉} (c3);
\draw[amplify, bend left=15] (c3) to node[below, font=\scriptsize, text=red!70] {离线错觉掩盖缺陷} (c4);
\draw[amplify, bend left=15] (c4) to node[left, font=\scriptsize, text=red!70] {评测失真延误识别} (c1);

% Center labels
\node[note, text width=5cm, below=3mm of base] {四类挑战形成相互放大循环};
\node[note, text width=4cm, above=12mm of c1, xshift=28mm] {单一挑战不可孤立解决};
\end{tikzpicture}
}
\caption{关键挑战的相互放大关系：四类结构性问题通过统一底座形成闭环反馈}
\label{fig:challenge-amplification}
\end{figure}

\subsection{执行与通信主路径：平台异构的持续扩散}

第 3 节已将后端能力契约列为执行路径的核心问题，本节关注的是：为什么平台异构在国产推理引擎中不是一次性适配问题，而是沿执行路径持续扩散的系统风险。

平台异构的特殊难点在于硬件代际能力差异大且缺少统一适配层。vLLM-Ascend、vLLM-MLU 与 SGLang Ascend 的公开材料都体现出对上游主线版本、CANN SDK 和 torch\_npu 三方联合对齐的强依赖\citep{vllmascend,vllmmlu,sglang_ascend_npu}。不同硬件上的 feature parity 难以默认成立——graph mode、attention backend、量化、PD 分离和 speculative decoding 等能力即使在同一工具链中也可能存在支持粒度差异\citep{fastdeploy_docs,lmdeploy_supported_models}。公开性能数字不能直接横向比较，因为模型版本、输入/输出长度、并发度、精度和 runtime 版本的任一变化都可能改变结论。

\subsection{状态治理主路径：复杂任务场景的能力补齐}

当前真实负载已从单轮文本生成扩展到长上下文、多模态、结构化输出、工具调用和多 Agent 协同\citep{dong2024xgrammar,liu2025elasticmm,xie2025aimetropolis}。推理引擎的优化目标不再是单一路径 decode，而是跨阶段状态管理：结构化输出要求把约束执行压进采样环路，工具调用要求支持生成与外部执行之间的往返切换，多模态请求要求消化更长的前处理链和更剧烈的 shape 波动。对国产平台，难点不在接口名义上"可用"，而在这些能力会同时冲击调度、缓存和图模式。

复杂任务场景直接改写了"状态"的性质。文本 serving 时代系统重点管理请求队列、KV Cache 和 batch 形状；进入多模态和工具调用场景后，运行时还须管理视觉前缀、语法状态、工具回注上下文、外部执行等待与流式中断恢复等状态对象。它们既非纯粹的模型内部变量，也无法完全交给服务层处理，而是直接影响调度顺序、缓存复用和回退路径选择的运行时实体。

\begin{figure}[!htb]
\centering
\resizebox{0.85\linewidth}{!}{\begin{tikzpicture}[
    node distance=5mm and 4mm,
    stage/.style={draw, rounded corners=2pt, minimum width=22mm, minimum height=10mm, align=center, font=\small, inner sep=4pt},
    side/.style={draw, rounded corners=2pt, align=center, font=\scriptsize, inner sep=3pt},
    flow/.style={-{Latex[length=2.2mm]}, semithick}
]
\node[stage, fill=gray!10] (ingress) {请求进入\\文本/多模态/工具调用};
\node[stage, fill=blue!10, right=of ingress] (normalize) {请求整形\\长度 bucket / SLO};
\node[stage, fill=green!10, right=of normalize] (prefill) {前处理与 Prefill\\视觉/语音编码};
\node[stage, fill=green!10, right=of prefill] (decode) {Decode\\采样/约束执行};
\node[stage, fill=orange!10, right=of decode] (verify) {验证与回退\\grammar / speculative};
\node[stage, fill=gray!10, right=of verify] (finish) {响应返回\\状态释放与观测};

\node[side, fill=yellow!8, below=9mm of prefill] (cache) {共享状态\\KV Cache / prefix / 迁移元数据};
\node[side, fill=yellow!8, below=9mm of decode] (tool) {外部交互\\工具执行 / 回注 / 再进入};

\draw[flow] (ingress) -- (normalize);
\draw[flow] (normalize) -- (prefill);
\draw[flow] (prefill) -- (decode);
\draw[flow] (decode) -- (verify);
\draw[flow] (verify) -- (finish);
\draw[flow, dashed] (prefill) -- (cache);
\draw[flow, dashed] (decode) -- (cache);
\draw[flow, dashed] (verify) -- (cache);
\draw[flow, dashed] (decode) -- (tool);
\draw[flow, dashed] (tool.west) |- (normalize.south);

\node[font=\scriptsize, align=center, below=8mm of normalize] {复杂请求的关键难点不在单步执行，而在跨阶段状态传递、资源复用与回退一致性。};
\end{tikzpicture}}
\caption{复杂请求生命周期示意图}
\label{fig:complex-request-lifecycle}
\end{figure}

推测解码进一步抬高了对执行一致性的要求。SpecInfer、Medusa、EAGLE 与 Lookahead Decoding 分别代表辅助模型、多头预测、特征不确定性建模与并行前瞻路线\citep{specinfer2023,leviathan2023fast,cai2024medusa,li2024eagle,fu2024lookahead}。硬件侧经验表明收益高度条件化：在 MI300X 上 speculative decoding 主要受益于小 batch 的带宽受限 decode，batch 增大后收益明显收缩\citep{singh2025specdecmi300x}。这类机制对国产平台尤其敏感，因为本地硬件在图模式切换和动态 shape 支持上约束更强——一个"只是多一次验证"的机制，迁移到新后端后可能同时意味着图缓存失效和异常回退链变长。

\subsection{压缩与成本主路径：收益不稳定与闭环缺失}

压缩路径的挑战与执行路径和状态治理直接耦合。执行路径上的平台异构会限制低比特算子的可用边界，状态治理层的缓存组织则决定压缩策略能否在运行时稳定兑现——三者缺一，压缩收益就难以从离线实验转化为生产能力。

"低成本"不是部署阶段最后再追问的经济性指标，而是与执行效率和状态治理同级的系统目标。问题在于压缩策略的收益往往不线性兑现：同一组量化参数在不同后端、不同上下文长度和不同负载强度下，可能分别表现为吞吐提升、尾时延恶化或精度波动。QServe 说明只有把量化、kernel 组织和 serving runtime 一起设计，低比特收益才可能在真实服务路径中落地\citep{wang2024qserve}；KIVI 与 SnapKV 表明长上下文场景下的状态压缩如果不和恢复路径与质量门限一起考虑，容量节省就会转化为恢复时延或命中率损失\citep{liu2024kivi,li2024snapkv}。对国产平台，这种耦合还会被放大——许多后端对低比特算子、动态 shape 和异步搬运的支持边界不一致。

\subsection{统一底座：上游演进与评测体系}

执行、状态和压缩三条主路径上的挑战，最终都需要回到“统一底座”层面来解决。所谓统一底座，不仅指技术架构上的抽象与观测基础设施，更指系统能否在持续演进中维持工程秩序——这涉及上游演进策略与评测体系两个维度。

高性能推理框架仍在快速演化。国产后端如果通过侵入式修改共享路径获得短期收益，后续将面临高昂合并成本。本地分叉之所以危险，还因为它逐步改变团队的优化激励——若长期依赖复制上游模块并在本地修补，团队最容易积累的是"快速打补丁"而非"设计稳定扩展面"的经验。相反，若平台特化更多收敛在插件、注册表和受控配置门面中，团队就更有机会把适配劳动沉淀为可复用的工程机制。因此 merge-safe 不是"对上游友好"的附属要求，而是国产推理引擎能否持续演进的核心条件。

评测体系方面，当前不少优化仍主要围绕静态 benchmark 展开。但真实业务包含请求长度分布变化大、多租户并发、冷热模型混部和 SLA 约束严格等特征。若评测只覆盖理想化场景，就难以反映系统在生产环境中的优势与短板。未来评测需要从"单点性能测试"扩展为"端到端工作负载评测"，将 TTFT、稳定吞吐、缓存命中率、异常恢复时间和资源成本纳入统一分析框架。OpenCompass 与 FlagPerf 分别代表了模型能力评测和系统性能评测的生态努力\citep{opencompass,flagperf}。真正有价值的评测应当连接这两类视角——既看模型效果，也看支撑这些效果所付出的系统成本。评测体系需要分层：底层算子用 microbenchmark 验证，运行时调度用受控 serving workload 验证（重点是 TTFT、TPOT 与尾时延），而平台演进能力和运维代价则更适合通过持续集成基线、版本升级回归和故障演练来评估。

\section{未来研究展望}

第 5 节将路线比较收束为若干系统挑战，本节进一步讨论未来更值得持续投入的研究方向。这些方向可压缩为三条主线：面向异构平台的能力契约与共享抽象、面向真实交付的控制面与证据链治理，以及面向复杂请求图的状态中心化运行时。图\ref{fig:outlook-map} 概括了三条主线在国产推理生态中的收敛关系。

\begin{figure}[!htb]
\centering
\resizebox{0.94\linewidth}{!}{\begin{tikzpicture}[
    node distance=18mm and 24mm,
    block/.style={draw=black!75, line width=0.8pt, rounded corners=2.2pt, minimum width=31mm, minimum height=10mm, align=center, font=\small, inner sep=4pt},
    centerblock/.style={draw=black!75, line width=0.8pt, rounded corners=3pt, minimum width=40mm, minimum height=12mm, align=center, font=\small\bfseries, inner sep=5pt},
    note/.style={font=\footnotesize, align=center, fill=white, inner sep=1.2pt, text=black!80},
    flow/.style={-{Latex[length=2.2mm]}, semithick, draw=black!75}
]
\node[centerblock, fill=gray!8] (goal) {可持续的国产算力\\推理系统底座};

\node[block, fill=blue!9, above left=20mm and 28mm of goal] (plugin) {后端插件化\\平台抽象/注册/可合并优化};
\node[block, fill=green!10, above right=20mm and 28mm of goal] (runtime) {资源管理与调度\\SLO/迁移/KV Cache/阶段解耦};
\node[block, fill=orange!12, below left=20mm and 28mm of goal] (eval) {统一评测与观测\\近真实负载/成本/稳定性};
\node[block, fill=red!8, below right=18mm and 24mm of goal] (model) {模型结构持续演进\\MoE/MLA/多模态/语音交互};

\draw[flow] (plugin.south east) -- (goal.north west);
\draw[flow] (runtime.south west) -- (goal.north east);
\draw[flow] (eval.north east) -- (goal.south west);
\draw[flow] (model.north west) -- (goal.south east);

\node[note, text width=6.2cm, above=6mm of goal] {未来竞争将从单点热点优化转向系统能力收敛。};
\node[note, text width=6.4cm, below=6mm of goal] {平台抽象、运行时、评测与模型演进需要同步推进。};
\end{tikzpicture}}
\caption{国产推理引擎未来演进方向关系图}
\label{fig:outlook-map}
\end{figure}

这三条主线并不是彼此独立的新功能列表，而是前文"一个底座、三条主路径"在未来更可能沉淀出的高层研究议题。它们之间存在递进关系：能力契约为系统演进提供稳定边界，控制面治理让这些边界可验证、可交付，而请求图重构则要求前两者同时升级以适应更复杂的运行时状态。

\subsection{面向异构平台的能力契约与共享抽象}

核心问题是：平台差异能否从"项目内部经验"上升为"生态可以共享的能力契约"。具体而言，需要回答以下开放研究问题：

\begin{enumerate}[leftmargin=*, label=(\arabic*)]
\item 后端能力应如何形式化表达，才能既覆盖 graph/eager 约束、通信语义和低比特路径的真实差异，又不把短期平台特例固化为长期接口？现有工程实践中，能力边界通常以隐式特判存在，缺少统一描述语言。
\item 回退语义应如何定义，才能保证图模式切换、缓存状态和分布式执行在回退前后保持一致？当前系统对回退的处理通常是"降级到 eager"，缺乏系统性的状态保持保证。
\item 共享抽象如何与平台特化协同演进？若抽象过薄，平台差异反复溢出；若过厚，又把短期特例固化。Llumnix、SOLA、LServe 等工作之所以重要，正因为它们显示运行时组织方式正在脱离"单一模型 serving 外壳"，转向资源与状态中枢的形态\citep{sun2024llumnix,hong2025sola,yang2025lserve}。
\end{enumerate}

只有当平台差异能够被收敛为共享能力矩阵而非反复改写系统主链时，新模型、新后端和新任务形态的进入才会从高风险迁移变成稳定的生态扩展。然而，能力契约本身并不能自动保证这些边界在生产环境中被持续遵守——这需要控制面来承担验证、回归和交付的闭环责任。

\subsection{控制面、评测与交付的一体化治理}

核心问题是：控制面能否从运维脚本集合演化为连接评测、成本与交付的一体化治理层。具体开放问题包括：

\begin{enumerate}[leftmargin=*, label=(\arabic*)]
\item 近真实工作负载应如何标准化表达？当前缺少能同时覆盖请求分布、上下文长度变化、多模态混合和工具调用比例的统一 workload 描述格式。BurstGPT 等工作已开始提供真实负载数据集\citep{wang2025burstgpt}，但从数据集到可复现评测协议之间仍有显著距离。
\item 版本基线、回退语义与故障注入应如何统一？升级回归和灰度成本目前无法被量化治理，大多停留在人工经验阶段。
\item 模型效果、系统成本与交付代价如何进入同一证据链？避免"性能结论"与"生产结论"长期脱节。
\end{enumerate}

执行面负责 token 生成，控制面负责让 token 生成处于可比较、可复现、可灰度和可回滚的条件之下。只有控制面被视为与执行面并列的一等对象，前文讨论的成本、评测与交付问题才可能从局部经验转化为稳定的生态能力。控制面治理的成熟度，反过来也取决于运行时本身能否提供更清晰的状态对象——这正是第三条主线关注的核心问题。

\subsection{面向多阶段请求图的运行时重构}

核心问题是：随着文本 decode、视觉 encode、语音前端、结构化约束验证、工具调用回注和多 Agent 协同被纳入同一服务链路，推理引擎如何从维护单一 token 生成循环转向管理多阶段请求图。具体开放问题包括：

\begin{enumerate}[leftmargin=*, label=(\arabic*)]
\item 请求图中的状态转移应如何显式建模？DeepSeek-V2/V3 的 MLA 与 MoE 组织连同 FlashMLA、Mooncake 等配套工作，已把通信组织、缓存布局和注意力实现重新推回运行时中心位置\citep{liu2024deepseekv2,deepseekv3,deepseekflashmla,qin2024mooncake}。这要求状态对象不再只按模型名单组织，而应围绕阶段关系组织。
\item 共享缓存与编译产物应如何围绕阶段关系而非模型名单组织？新任务形态进入时不应需要重写主链。
\item 运行时如何同时容纳模型 co-design、状态治理和交付控制面？Qwen2-VL、InternVL2、GLM-4-Voice 与 MiniCPM-o 等模型已说明视觉、视频和语音链路使文本时代的单一路径假设难以维持\citep{wang2024qwen2vl,chen2024internvl2,zeng2024glm4voice,cui2026minicpmo45}。
\end{enumerate}

这个转变未必立刻带来最亮眼的跑分，但更可能决定国产推理系统能否在未来几年内持续吸收新结构、新任务和新平台。系统边界会被模型公司、应用形态和硬件路线共同牵引，任何一方的变化都会迫使另外两方重写分工。

综合来看，这三条主线共同指向一个判断：国产推理引擎的下一阶段竞争，将从"谁更快适配某个硬件后端"转向"谁能在异构平台、复杂负载和快速演进之间建立更稳定的工程秩序"。能力契约、控制面治理和请求图重构不是三个孤立议题，而是同一问题的三个切面——系统能否把持续增加的复杂度组织进可扩展的边界，而不是让它们在外围不断堆积。

\section{结束语}

本文围绕"生态格局—分析骨架—路线比较—关键挑战—收敛方向"这一主线，对国产算力平台上的大模型推理引擎作了结构化梳理。全文的核心判断可归纳为：国产推理引擎的竞争点不在局部算子是否足够快，而在于系统能否在同一架构中同时承受后端异构、复杂任务负载和上游快速演进三类压力，并在国产芯片与本土部署环境中形成可持续的工程闭环。

具体而言，这一判断体现在四个方面。第一，平台异构的真正终点不是兼容，而是硬件差异能否被收敛为稳定、可验证、可回退的能力契约——若能力边界长期以隐式特判存在，新后端的接入就会反复冲击共享主链。第二，长上下文、多模态、结构化输出与 Agent 负载正在把推理引擎从 token 生成器推向状态中心化运行时，状态治理能力将比局部算子优化更能决定系统上限。第三，量化与 KV 压缩只有进入统一的观测、调度和成本闭环，才可能从离线实验结果转化为可持续交付的单位 token 成本优势。第四，路线选择本质上是复杂度分配与责任边界选择——复杂度由 backend、runtime、state machine 还是控制面承担，将直接决定系统的演进速度与长期维护成本。

面向未来，国产推理引擎的下一阶段竞争不在单点峰值还能提高多少，而在于谁能更快、更稳、更低成本地吸收模型变化并完成生产交付。对产业侧而言，建设重点应从围绕单模型做专项优化转向围绕持续演进负载建设系统能力；对研究侧而言，则需要把平台迁移成本、长周期维护成本和近真实服务指标放到与性能同等重要的位置。围绕这些问题持续推进，国产推理引擎的演进才更可能从阶段性适配走向可持续积累。


% ===========================================
% Author bio (CCCF style: before references)
% ===========================================
\vspace{1em}
\noindent\begin{minipage}{\linewidth}
\noindent % Author photo placeholder
\textbf{金海}~~CCF会士、副理事长。华中科技大学计算机学院教授。主要研究方向为计算机体系结构、计算系统虚拟化、集群计算和云计算、并行I/O。hjin@hust.edu.cn
\end{minipage}

\vspace{1.2em}

% ===========================================
% References
% ===========================================
\begingroup
\sloppy
\bibliographystyle{unsrtnat}
\bibliography{refs}
\endgroup

% ===========================================
% English title, abstract & keywords (CCCF style: at end)
% ===========================================
\vspace{1.5em}
\noindent\textbf{\large A Survey of Large-Model Inference Engines on Domestic Computing Platforms}

\vspace{0.6em}
\noindent\textbf{Zhang Shuhao, Liu Shifeng, Du Yufeng, Ma Junhao, Qiu Ruijie, Liao Xiaofei, Jin Hai}

\vspace{0.6em}
\noindent\textbf{Abstract}\quad This survey examines large-model inference engines on domestic computing platforms. Using cross-checkable public evidence, it explains why the ecosystem has split into three routes---upstream-core reuse and enhancement, vertically integrated platform stacks, and domestic-native systems---and adopts an analytical scaffold comprising a unified abstraction/observability base plus three main paths (execution and communication, state governance, compression-cost co-optimization) to compare where each route places its primary complexity. The survey distills structural challenges around hardware heterogeneity, complex workloads, unstable compression gains, and rapid upstream evolution, and outlines future directions including platform capability contracts, state-centric runtime restructuring, and unified control-plane governance.

\vspace{0.4em}
\noindent\textbf{Key Words}\quad domestic computing platforms; large-model inference; inference engine; runtime systems; hardware-software co-design

% ===========================================
% Chinese abstract & keywords (CCCF style: at end)
% ===========================================
\vspace{1.2em}
\noindent\textbf{摘要}\quad 本文面向国产算力平台上的大模型推理引擎，围绕"生态格局—分析骨架—路线比较—关键挑战—未来收敛"这一主线展开结构化综述。以可交叉核验的公开证据为基础，本文首先梳理当前生态为何分化为开源主干复用增强、平台一体化方案与国产原生/独立路线三类，继而以"一个底座（统一抽象与可观测证据链）、三条主路径（执行与通信、状态治理、压缩与成本协同）"为分析骨架，比较不同路线把复杂度主要压在何处。在此基础上，将平台异构、复杂任务负载、压缩收益不稳定、上游快速演进与评测体系失真等问题收敛为结构性挑战，并讨论能力契约沉淀、状态中心化运行时重构和控制面一体化治理等未来方向，以期为国产推理引擎的路线选择与架构演进提供可复核的分析框架。

\vspace{0.4em}
\noindent\textbf{关键词}\quad 国产算力平台；大模型推理；推理引擎；运行时；软硬协同

\vspace{0.4em}
\noindent\textbf{中图分类号：}TP39

\end{spacing}

\end{document}
